\documentclass[11pt]{article}

\usepackage[margin=1in]{geometry}
\usepackage{amsmath,amssymb,amsthm}
\usepackage{booktabs}
\usepackage{enumitem}
\usepackage[hidelinks]{hyperref}

\newtheorem{observation}{Observation}
\newtheorem*{question}{Open question}
\theoremstyle{definition}
\newtheorem{remark}{Remark}

\newcommand{\No}{\mathbf{No}}
\newcommand{\On}{\mathrm{On}_2}
\newcommand{\F}{\mathbb{F}}
\newcommand{\Tr}{\operatorname{Tr}}
\newcommand{\Arf}{\operatorname{Arf}}
\newcommand{\nimmul}{\otimes}
\newcommand{\Cl}{\mathrm{Cl}}

\title{Clifford algebras over the field-cores of combinatorial games,\\
and the Arf invariant as a win-bias}
\author{a9lim}
\date{June 2026}

\begin{document}
\maketitle

\begin{abstract}
Conway's combinatorial games form a partially ordered abelian group under
disjunctive sum but not a ring, so a Clifford algebra---which needs a
commutative scalar ring---can only be built over the field-like subclasses of
game-world: the surreal numbers $\No$, the surcomplex numbers $\No[i]$, and the
nimbers $\On$. We describe \texttt{pleroma}, a small library realising Clifford
algebras (with nilpotents) over all three, and focus on the characteristic-$2$
nimber case, where the Clifford algebra is classified by the \emph{Arf
invariant} of an $\F_2$ quadratic form (Ovsienko). This is an experimental note.
Its modest contributions are: (i) an implementation that keeps the quadratic and
polar forms independent, so the char-$2$ theory is faithful, and that computes
the Arf invariant over any nim-subfield via the field trace; (ii) the
observation, verified computationally, that the Arf-bearing Gold forms
$\Tr(x^{1+2^a})$ are \emph{composites of combinatorial-game operations}; and
(iii) the framing of the Arf invariant as a \emph{win-bias}: by Dickson's
democratic-value theorem it is exactly the sign of the majority margin of the
form's zero-set, which we confirm on the game-built forms. We close with a sharp
open question: is there a \emph{natural} game whose $P$-positions are such a
form's zero set? The classical results we lean on (the Arf--Clifford
correspondence, Dickson's theorem, the Gold-function ranks, and
Turning-Corners~$=$~nim-multiplication) are reproduced by the tool as
validation, not claimed as new.
\end{abstract}

\section{Why three backends}

A (short, partizan) combinatorial game is a pair $\{L \mid R\}$ of sets of
previously-constructed games; under the disjunctive sum the games form a
partially ordered abelian group~\cite{Conway76,WW}. They do \emph{not} form a
ring: Conway's game product is only a congruence on the \emph{numbers}, the
surreal subclass~$\No$. A Clifford algebra is a quotient of a tensor algebra and
requires a commutative ring of scalars, so ``Clifford over games'' is forced
onto the field-like cores of game-world:

\begin{center}
\begin{tabular}{lll}
\toprule
core & structure & Clifford flavour \\
\midrule
$\No$ (surreals) & real-closed field, char $0$ & the $\Cl(p,q)$ theory over $\mathbb{R}$, exotic scalars \\
$\No[i]$ (surcomplex) & algebraically closed, char $0$ & the complex Clifford theory \\
$\On$ (nimbers) & algebraically closed, char $\mathbf{2}$ & genuinely different \\
\bottomrule
\end{tabular}
\end{center}

\noindent
$\No$ is real-closed~\cite{Conway76}, so finite-dimensional Clifford algebras
over it are classified exactly as over $\mathbb{R}$ (the $8$-fold periodicity),
but one may now use infinite or infinitesimal metric entries such as
$e_0^2=\omega$, $e_1^2=\varepsilon=\omega^{-1}$. $\No[i]$ is algebraically
closed and gives the complex theory. The nimbers $\On$---the ordinals under
nim-addition and nim-multiplication---form an algebraically closed field of
\emph{characteristic $2$}~\cite{Conway76}, with the finite subfields
$\F_{2^{2^k}}$ consisting of the nimbers below $2^{2^k}$. This is the only core
where Clifford theory acquires a new character, and it is our focus.

\section{Characteristic $2$, faithfully}

In characteristic $2$ a quadratic form is not determined by a symmetric bilinear
form. Writing $q_i = Q(e_i)$ for the squares and $b_{ij}=Q(e_i+e_j)+Q(e_i)+Q(e_j)$
for the polar form, the relations defining the algebra are
\[
  e_i^2 = q_i, \qquad e_i e_j + e_j e_i = b_{ij},
\]
and the polar form is \emph{alternating} ($b_{ii}=0$) while $q_i$ may be nonzero.
If one collapses $(q,b)$ into a single symmetric bilinear form, every char-$2$
algebra over an orthogonal basis becomes commutative and the content is lost.
\texttt{pleroma} therefore stores $q$ and $b$ as independent data and reduces
products by a word-rewriting routine whose only use of the sign is through the
scalar's own negation; since $-1=1$ in $\On$, the char-$2$ sign-vanishing falls
out automatically.

For a nonsingular quadratic form $Q$ on $\F_2^{2m}$ with symplectic basis
$\{a_k,b_k\}$, the \emph{Arf invariant} is
\(
  \Arf(Q) = \sum_{k} Q(a_k)\,Q(b_k) \in \F_2
\)~\cite{Arf41};
it is independent of the basis and is a complete invariant. The crucial fact for
us is the classification it controls.

\begin{quotation}
\noindent\textbf{Theorem (Ovsienko~\cite{Ovsienko16}).} The classification of
(real) Clifford algebras is equivalent to the classification of quadratic forms
over $\F_2$, with the Arf invariant the complete invariant.
\end{quotation}

\noindent
So computing $\Arf$ of a nim-Clifford form returns the isomorphism class of the
algebra: it is the char-$2$ Clifford classifier, not a toy. \texttt{pleroma}
computes it for a form over any nim-subfield $\F_{2^{2^k}}$ by symplectic
reduction over the field, then pushes the result to $\F_2$ via the trace
$\Tr_{\F_{2^m}/\F_2}(x)=x+x^2+\cdots+x^{2^{m-1}}$, which realises the canonical
isomorphism $k/\wp(k)\cong\F_2$. Over $\F_4$, for instance, the form with
$q=(\ast 2,\ast 2)$ is of type $O^-$ and $q=(\ast 2,\ast 3)$ is hyperbolic.

\section{The bridge to games}

Nim-addition is the disjunctive sum (Sprague--Grundy), and nim-multiplication is
itself a game. Conway's \emph{Turning Corners} has, for the single coin at
$(x,y)$, the Grundy value given by the excludant recurrence
\begin{equation}\label{eq:tc}
  x \nimmul y \;=\; \operatorname{mex}\bigl\{\,(i\nimmul y)\oplus(x\nimmul j)\oplus(i\nimmul j)
        \;:\; 0\le i<x,\ 0\le j<y \,\bigr\},
\end{equation}
which is Conway's definition of nim-multiplication~\cite{Conway76,WW}. We
implement~\eqref{eq:tc} directly and verify it agrees with the algebraic
Fermat-power product (computed via $F_n\nimmul F_n=\tfrac32 F_n$ for Fermat
$2$-powers $F_n=2^{2^n}$) on all pairs $x,y<48$: two independent definitions,
one combinatorial and one field-theoretic, of the same product.

This makes three operations game-realizable: the nim-product $\nimmul$ (Turning
Corners), the Frobenius $\square\colon v\mapsto v^2 = v\nimmul v$ (the diagonal
of Turning Corners), and $\oplus$ (disjunctive sum); hence also the trace, which
is iterated $\square$ and $\oplus$.

\paragraph{Gold forms, and a validation.}
The natural quadratic forms on a char-$2$ field are
$Q_a(x)=\Tr\!\left(x\cdot x^{2^a}\right)=\Tr\!\left(x^{1+2^a}\right)$, the
\emph{Gold forms}~\cite{Gold68}: $g=\mathrm{Frob}^a$ is additive, so $Q_a$ is
genuinely quadratic. As a standard cross-check, the polar-form rank our
classifier computes reproduces the known Gold-function rank
$\operatorname{rank}=m-2\gcd(a,m)$ (in the regime $m=2^k$, where $m/\gcd$ is
even)~\cite{LN}; see Table~\ref{tab:gold}.

\begin{table}[h]
\centering
\begin{tabular}{rrccccc}
\toprule
field & $a$ & $Q$ & $\Arf$ & type & rank & $m-2\gcd(a,m)$ \\
\midrule
$\F_{2^{4}}$  & $1$ & $\Tr(x^3)$    & $1$ & $O^-$ & $2$  & $2$ \\
$\F_{2^{8}}$  & $1$ & $\Tr(x^3)$    & $1$ & $O^-$ & $6$  & $6$ \\
$\F_{2^{8}}$  & $2$ & $\Tr(x^5)$    & $1$ & $O^-$ & $4$  & $4$ \\
$\F_{2^{16}}$ & $1$ & $\Tr(x^3)$    & $1$ & $O^-$ & $14$ & $14$ \\
$\F_{2^{16}}$ & $4$ & $\Tr(x^{17})$ & $1$ & $O^-$ & $8$  & $8$ \\
$\F_{2^{32}}$ & $1$ & $\Tr(x^3)$    & $1$ & $O^-$ & $30$ & $30$ \\
\bottomrule
\end{tabular}
\caption{The classifier's rank matches the Gold-function formula
(15/15 cases up to $m=32$; a representative selection).}
\label{tab:gold}
\end{table}

\paragraph{The Gold forms are game-built.}
Because $\nimmul$, $\square$, $\oplus$, and the trace are all game operations,
the Gold form
\(
  Q_a(v) = \Tr\!\bigl(v \nimmul v^{2^a}\bigr)
\)
is a composite of combinatorial-game operations applied to a position's nimber
value $v$ (and under the $1$-dimensional coin-turning game with single-coin
Grundy value $g(n)=2^n$, a position's value \emph{is} a nimber). We confirm this
by rebuilding $Q_a$ from literal Turning-Corners products and checking equality
with the field computation on the small fields where the
recurrence~\eqref{eq:tc} is tractable.

\begin{observation}\label{obs:built}
The quadratic forms whose Arf invariants classify the nim-Clifford algebras are
not merely defined over the value field of games; they are composites of game
operations. The char-$2$ Clifford theory is, in this constructive sense, made of
games.
\end{observation}

\section{The Arf invariant as a win-bias}

The semantic payoff comes from a classical counting fact.

\begin{quotation}
\noindent\textbf{Theorem (Dickson~\cite{Dickson01}; Arf~\cite{Arf41}).}
Over $\F_2$, the Arf invariant of a quadratic form is the value the form takes
most often. For a nonsingular form on $\F_2^{2m'}$,
\[
  \#\{\,v : Q(v)=0\,\} \;=\; 2^{2m'-1} + (-1)^{\Arf(Q)}\,2^{m'-1}.
\]
\end{quotation}

\noindent
Read game-theoretically: if a game had $P$-positions (second-player wins)
exactly $\{v:Q(v)=0\}$, then $\Arf(Q)$ would be precisely the \emph{sign of the
win-bias}---which player wins from more starting positions---with a fixed margin
of $2^{m'-1}$, a Gauss-sum, i.e.\ a square-root-scale fluctuation about
$50/50$. The Arf invariant is the second-order tiebreaker of an otherwise
near-fair game.

We verify the count directly on the game-built Gold forms, handling the radical:
for rank $r$ and an isotropic radical, $\#\{Q=0\}=2^{m-1}+(-1)^{\Arf}2^{m-r/2-1}$
(and exactly $2^{m-1}$ when the radical is anisotropic). The brute-forced
zero-counts match in every case (Table~\ref{tab:bias}).

\begin{table}[h]
\centering
\begin{tabular}{rrcccr}
\toprule
field & $a$ & $\Arf$ & rank & $\#\{Q=0\}$ & bias \\
\midrule
$\F_{2^{4}}$  & $1$ & $1$ & $2$  & $4$     & $-4$ \\
$\F_{2^{8}}$  & $1$ & $1$ & $6$  & $112$   & $-16$ \\
$\F_{2^{8}}$  & $2$ & $1$ & $4$  & $96$    & $-32$ \\
$\F_{2^{16}}$ & $1$ & $1$ & $14$ & $32512$ & $-256$ \\
$\F_{2^{16}}$ & $4$ & $1$ & $8$  & $30720$ & $-2048$ \\
\bottomrule
\end{tabular}
\caption{Zero-counts of the game-built Gold forms match
$2^{m-1}+(-1)^{\Arf}2^{m-r/2-1}$ (bias $=\#\{Q=0\}-2^{m-1}$);
10/10 cases verified.}
\label{tab:bias}
\end{table}

\begin{observation}
On the game-built forms, the Arf invariant equals the win-bias in the counting
sense: it is the sign of the majority margin of the zero-set.
\end{observation}

\section{What remains open}

The bridge closes at the level of \emph{construction} (Observation~\ref{obs:built})
and of \emph{counting} (the Arf invariant is the zero-set bias). What is not
established is a \emph{realisation}: a natural game whose $P$-positions are
$\{v:Q(v)=0\}$. The obstruction is structural and informative.

\begin{remark}
Normal-play disjunctive sums cannot do it. Their outcomes are $\oplus$-linear in
the components (Sprague--Grundy), so a $P$-position set is always a linear
subspace, never a genuine quadric. This is exactly why coin-turning yields the
\emph{bilinear} nim-product but not a quadratic refinement.
\end{remark}

\begin{question}
Is there a natural, necessarily non-disjunctive game whose $P$-positions are the
zero set of a Gold form $Q_a$ (equivalently, whose win-bias is read off by
$\Arf(Q_a)$)? Two candidate families avoid the linearity obstruction:
\emph{interactive} coin-turning / lights-out games coupled through the polar form
$b$, and \emph{misère} play, where sums are genuinely non-linear and quadratic
structure can live in a misère quotient~\cite{WW}. We are not aware of either in
the literature; Welter's game, the deepest relative, lands in the representation
theory of symmetric groups~\cite{Irie18} rather than in Arf invariants.
\end{question}

A candidate must hit the win-bias of Table~\ref{tab:bias}; \texttt{pleroma}
makes that target computable. The source, tests, and the three experiments
underlying Tables~\ref{tab:gold} and~\ref{tab:bias} are available at
\url{https://github.com/a9lim/pleroma}.

\section*{Acknowledgements}
The \texttt{pleroma} library, the experiments, and this note were developed in
collaboration with Claude (Anthropic), which wrote the Rust/Python
implementation and ran the computational probes; direction, framing, and
verification are the author's.

\begin{thebibliography}{9}
\bibitem{Arf41} C.~Arf, \emph{Untersuchungen \"uber quadratische Formen in
K\"orpern der Charakteristik~2, I}, J.\ Reine Angew.\ Math.\ \textbf{183}
(1941), 148--167.
\bibitem{WW} E.~R.~Berlekamp, J.~H.~Conway, R.~K.~Guy, \emph{Winning Ways for
Your Mathematical Plays}, 2nd ed., A~K Peters, 2001--2004.
\bibitem{Conway76} J.~H.~Conway, \emph{On Numbers and Games}, Academic Press,
1976.
\bibitem{Dickson01} L.~E.~Dickson, \emph{Linear Groups, with an Exposition of
the Galois Field Theory}, Teubner, 1901.
\bibitem{Gold68} R.~Gold, \emph{Maximal recursive sequences with $3$-valued
recursive cross-correlation functions}, IEEE Trans.\ Inform.\ Theory \textbf{14}
(1968), 154--156.
\bibitem{Irie18} Y.~Irie, \emph{$p$-Saturations of Welter's game and the
irreducible representations of symmetric groups}, J.\ Algebraic Combin.\
(2018); arXiv:1604.07214.
\bibitem{LN} R.~Lidl, H.~Niederreiter, \emph{Finite Fields}, 2nd ed., Cambridge
University Press, 1997.
\bibitem{Ovsienko16} V.~Ovsienko, \emph{Real Clifford algebras and quadratic
forms over $\F_2$: two old problems become one}, The Mathematical Intelligencer
\textbf{38} (2016); arXiv:1601.07664.
\end{thebibliography}

\end{document}
